\chapter{Введение: масштаб как история роста геометрии}
\label{ch:v10-introduction}

Том IX завершился запретом на происхождение абсолютного статического масштаба
из относительных спектральных, гауссовых и KMS-данных. Том X меняет предмет
поиска. Вместо ещё одного кандидата на постоянную энергию рассматривается
геометрическая история, в которой пространство расширяется посредством
добавления квантовых ячеек, а локальная физика задаётся устойчивым
безразмерным спектральным атласом.

Основная переменная нового тома --- число ячеек $N$. Относительный масштаб
пространства определяется формулой
\begin{equation}
 a=\left(\frac{N}{N_0}\right)^{1/3},
 \qquad
 \zeta=\log a=\frac13\log\frac{N}{N_0}.
 \label{eq:v10-intro-geometric-scale}
\end{equation}
Координата $\zeta$ одновременно является геометрическим параметром
ренормгруппового хода. Она не требует заранее выбранной единицы энергии и не
объявляется внешним физическим временем.

\section{Рабочая гипотеза}

Безразмерный спектральный оператор $\widehat D$ задаёт рисунок локальных
масс и связей, тогда как растущая история задаёт их размещение и
космологический ход:
\begin{equation}
 D(\zeta)=E_{\rm cell}(\zeta)\widehat D(g(\zeta)),
 \qquad
 \beta_i=\frac{dg_i}{d\zeta}.
 \label{eq:v10-intro-spectral-rg-factorization}
\end{equation}
Ветвь добавления ячеек сохраняет локальный размер ячейки. Ветвь совместного
растяжения всех локальных единиц считается лишь конформным переописанием до
появления изменяющегося безразмерного отношения.

\section{Контракт Тома X}

Положительное замыкание требует:
\begin{enumerate}
\item общего носителя истории роста, локального спектрального атласа и
квантовых полей;
\item происхождения нормированной меры переходов $N\to N+1$;
\item ненулевого квантового хода или следовой аномалии относительно $\zeta$;
\item связи темпа роста с космологической кривизной без подстановки
наблюдаемого значения;
\item типизированного отображения геометрического хода в локальный
энергетический сектор;
\item хотя бы одного безусловного безразмерного следствия.
\end{enumerate}

Первый гейт проверяет только общий носитель и точные кинематические тождества.
Физическое происхождение роста не предполагается заранее.

\section{Ритм интуитивного насыщения}

Локальная гейт-цепочка не должна расти быстрее понимания своего общего
родителя. После каждых трёх--пяти новых гейтов обязателен synthesis
checkpoint, включающий повторное чтение ключевых формул, сопоставление с
остальными ветвями программы и проверку первичной литературы. Следующий
локальный гейт разрешён только после фиксации того, какие механизмы были
действительно выведены, какие лишь условно сконструированы и какой новый
расчёт различает оставшиеся маршруты.

\begin{nogocriterion}[Не более пяти гейтов без интуитивного насыщения]
Последовательность из шести локальных гейтов без отдельной интуитивной
сводки считается методологически незамкнутой независимо от числа
полученных точных рангов и тождеств.
\end{nogocriterion}

\begin{protocol}[Открытие Тома X]\begin{enumerate}
\item вход: \textbf{масштабный запрет Тома IX};
\item новая переменная: \textbf{число геометрических ячеек $N$};
\item ренормгрупповая координата: \textbf{$\zeta=\frac13\log(N/N_0)$};
\item статическая калибровка: \textbf{не используется};
\item первый шаг: \textbf{допуск общего растущего носителя}.
\item cadence: \textbf{интуитивное насыщение после каждых $3$--$5$ гейтов}.
\end{enumerate}\end{protocol}